% Part: computability
% Chapter: recursive-functions
% Section: partial-functions

\documentclass[../../../include/open-logic-section]{subfiles}

\begin{document}

\olfileid[pt]{cmp}{rec}{par}
\olsection{Funções recursivas parciais}


Para motivar a definição das funções recursivas, observe que nossa
demonstração de que existem funções computáveis que não são recursivas
primitivas estabelece, na verdade, muito mais. O argumento foi simples:
tudo o que usamos foi o fato de ser possível enumerar funções
$f_0,f_1,\dots$ de tal modo que, como função de $x$ e $y$, $f_x(y)$
seja computável. Portanto, o argumento se aplica a \emph{qualquer classe
  de funções que possa ser enumerada dessa maneira}. Isso nos deixa em
um impasse: gostaríamos de descrever explicitamente as funções
computáveis; mas nenhuma descrição explícita de uma coleção de funções
computáveis pode ser exaustiva!

A saída é permitir que funções \emph{parciais} entrem em cena. Veremos
que \emph{é} possível enumerar as funções computáveis
parciais.\iftag{TMs}{ Na verdade, praticamente já sabemos que esse é o
  caso, pois é possível enumerar máquinas de Turing de maneira
  sistemática.}{} Voltaremos mais tarde a nosso argumento diagonal e
investigaremos por que ele não funciona quando funções parciais são
incluídas.

A pergunta agora é a seguinte: o que precisamos acrescentar às funções
recursivas primitivas para obter todas as funções recursivas parciais?
Precisamos fazer duas coisas:
\begin{enumerate}
\item Modificar nossa definição das funções recursivas primitivas para
  permitir também funções parciais.
\item \emph{Acrescentar} algo à definição, para que algumas novas funções
  parciais sejam incluídas.
\end{enumerate}

A primeira tarefa é fácil. Como antes, começaremos com zero, sucessor e
projeções, e fecharemos o conjunto sob composição e recursão primitiva.
A única diferença é que temos de modificar as definições de composição
e recursão primitiva para admitir a possibilidade de que alguns dos
termos na definição não estejam definidos. Se $f$ e $g$ são funções
parciais, escreveremos $f(x) \fdefined$ para indicar que $f$ está
definida em $x$, O.s.\@, que $x$ pertence ao domínio de $f$; e
$f(x) \fundefined$ para indicar o contrário, O.s.\@, que $f$ não está
definida em~$x$. Usaremos $f(x) \simeq g(x)$ para indicar que ou
$f(x)$ e $g(x)$ são ambas indefinidas, ou ambas estão definidas e são
iguais. Usaremos essas notações também para termos mais complicados.
Adotaremos a convenção de que, se $h$ e $g_0$, \dots,~$g_k$ são todas
funções parciais, então
\[
h(g_0(\vec x),\dots,g_k(\vec x))
\]
está definida se, e somente se, cada $g_i$ está definida em $\vec x$ e
$h$ está definida em $g_0(\vec x)$, \dots,~$g_k(\vec x)$. Com esse
entendimento, as definições de composição e recursão primitiva para
funções parciais são exatamente as anteriores, exceto por termos de
substituir ``$=$'' por ``$\simeq$''.

O que acrescentaremos à definição das funções recursivas primitivas
para obter funções parciais é o \emph{operador de busca ilimitada}. Se
$f(x,\vec z)$ é uma função parcial qualquer sobre os números naturais,
defina $\mu x \; f(x,\vec z)$ como
\begin{quote}
  o menor $x$ tal que $f(0,\vec z), f(1,\vec z), \dots, f(x,\vec
  z)$ estejam todos definidos e $f(x,\vec z) = 0$, se tal $x$ existir
\end{quote}
entendendo-se que $\mu x \; f(x,\vec z)$ é indefinido caso contrário.
Isso define $\mu x \; f(x,\vec z)$ de maneira única.

\begin{explain}
Observe que nossa definição não faz referência a\iftag{TMs}{ máquinas
  de Turing, nem a}{} algoritmos ou a qualquer modelo computacional
específico. Mas, assim como ocorre com a composição e a recursão
primitiva, existe uma intuição operacional e computacional por trás da
busca ilimitada. Quando se trata da computabilidade de uma função
parcial, os argumentos nos quais a função é indefinida correspondem a
entradas para as quais a computação não para. O procedimento para
calcular $\mu x \; f(x,\vec z)$ será o seguinte: calcule
$f(0,\vec z), f(1,\vec z), f(2,\vec z)$ até que seja devolvido o valor
0. Se, porém, algum dos cálculos intermediários não parar, o cálculo de
$\mu x \; f(x,\vec z)$ tampouco parará.
\end{explain}

Se $R(x,\vec z)$ é uma relação qualquer, define-se
$\mu x \; R(x,\vec z)$ como $\mu x \; (1 \tsub \Char{R}(x,\vec z))$.
Em outras palavras, $\mu x \; R(x,\vec z)$ devolve o menor valor de $x$
tal que $R(x,\vec z)$ vale. Assim, se $f(x,\vec z)$ é uma função total,
$\mu x \; f(x,\vec z)$ é o mesmo que $\mu x \; (f(x,\vec z) = 0)$.
Observe, contudo, que nossa definição original é mais geral, pois
admite a possibilidade de que $f(x,\vec z)$ não esteja definida em
todos os pontos (ao passo que, em contraste, a função característica
de uma relação é sempre total).

\begin{defn}
O conjunto das \emph{funções recursivas parciais} é o menor conjunto
de funções parciais dos números naturais nos números naturais (de
diversas aridades) que contém zero, sucessor e projeções e é fechado sob
composição, recursão primitiva e busca ilimitada.
\end{defn}

Naturalmente, algumas das funções recursivas parciais serão totais,
O.s.\@, definidas para todo argumento.

\begin{defn}
\ollabel{defn:recursive-fn}
O conjunto das \emph{funções recursivas} é
o conjunto das funções recursivas parciais que são totais.
\end{defn}

Uma função recursiva é às vezes chamada de ``recursiva total'' para
enfatizar que está definida em toda parte.

\end{document}


